\documentclass[../../main.tex]{subfiles}
\begin{document}

\begin{flushright}
  \footnotesize\textit{【自動文字起こし・要確認】}
\end{flushright}

{\bf [解]} (1) 1に4色すべてが入っている確率を $\alpha$ とすると，対称性から，求める確率 $P_1$ は
\begin{align}
  P_1 = \alpha^2 \cdots (1)
\end{align}
である．$\alpha$ について，5回の操作で，4つの色のうち，3つの色が1回ずつ，の残りの1色が2回出るので，
\begin{align}
  \alpha &= 4 \cdot {}_5C_2 \cdot {}_3C_1 \cdot {}_2C_1 \cdot \left(\dfrac{1}{4}\right)^5 \\
  &= \dfrac{15}{4^3}
\end{align}
だから，(1)より
\begin{align}
  P_1 = \dfrac{3^2 \cdot 5^2}{4^6} \quad \text{\#\!}
\end{align}

\paragraph{(2)}

5回の操作で，4色全てが少なくとも1回出れば良く，
\begin{align}
  P_2 = \alpha = \dfrac{3 \cdot 5}{4^3} \quad \text{\#\!}
\end{align}

\paragraph{(3)}

10回の操作で，4色全てが少なくとも2回出れば良い．この時，残りの2つの玉について，
\begin{itemize}
  \item[$1^\circ$] 2つとも同じ色
  \item[$2^\circ$] 2つとも違う色
\end{itemize}
がありうる．$1^\circ$ の時の確率は
\begin{align}
  4 \times {}_{10}C_{4} \cdot {}_{6}C_{2} \cdot {}_{4}C_{2} \cdot \left(\dfrac{1}{4}\right)^{10} = 3^3 \cdot 5^2 \cdot 7 \left(\dfrac{1}{4}\right)^8
\end{align}
$2^\circ$ の時の確率は
\begin{align}
  {}_{4}C_{2} \cdot {}_{10}C_{3} \cdot {}_{7}C_{3} \cdot {}_{4}C_{2} \cdot \left(\dfrac{1}{4}\right)^{10} = 3^3 \cdot 5^2 \cdot 7 \cdot 2 \left(\dfrac{1}{4}\right)^8
\end{align}
だから，
\begin{align}
  P_3 &= 3^3 \cdot 5^2 \cdot 7 \cdot \left(\dfrac{1}{4}\right)^8 (1+2) \\
  &= 3^4 \cdot 5^2 \cdot 7 \left(\dfrac{1}{4}\right)^8
\end{align}
(1)とあわせて，
\begin{align}
  \dfrac{P_3}{P_1} &= \dfrac{4^6 \cdot 3^4 \cdot 5^2 \cdot 7}{3^2 \cdot 5^2 \cdot 4^8} \\
  &= \dfrac{3^2 \cdot 7}{4^2} = \dfrac{63}{16} \quad \text{\#\!}
\end{align}

\end{document}
